%% ============================================================
%% CHAPTER 8 — Guarded Speculative Execution
%%                                        (budget: ~9,000 words)
%% Role: answers RQ4. Source: docs/specclaw-research-proposal.md.
%% Scope is adjustable: if time is short, cut §8.6.3 (adaptive depth)
%% and shrink the stochasticity grid; the chapter stands on RQ-a/b alone.
%% ============================================================
\chapter{Guarded Speculative Execution}
\label{ch:speculation}

\section{Motivation: the Fixed Cost of Every Call}
\label{sec:spec-motivation}
% The per-call fixed overhead (radio wake, TLS handshake, RTT) measured in
% Ch5; an n-step ReAct task pays it n times. Disconnection halts reasoning
% even though the loop survives. One paragraph recalling that prompt
% caching does not touch these costs.

\section{The Three Inversions}
\label{sec:spec-inversions}
% The table contrasting server-side speculation with this setting:
% verifier (target LLM -> physical world + deterministic guards);
% objective (latency overlap -> call count/energy/offline resilience);
% action semantics (discardable -> irreversible, commit-after-verify).
% Derive the new research object: silent divergence.

\section{The Speculation Script Protocol}
\label{sec:spec-protocol}
% The JSON script format ((action, guard envelope, timeout) tuples,
% irreversible flag, done_when). The minimal guard language (gt/lt/in/eq/
% delta_lt/within_s, AND only) and why restricted expressiveness is the
% point: decidable, ~100 LoC evaluator, statically auditable.
% Grammar in Appendix~\ref{app:guards}.

\section{Guarded-Commit Semantics}
\label{sec:spec-semantics}
% The execution algorithm (pre-guard -> D_irrev valve -> execute through
% the policy allowlist -> post-guard within timeout -> commit log).
% Definitions: acceptance length L; silent divergence; guard false-reject;
% D_irrev. The two invariants. The compact-state repredict payload
% (O(1) uplink per replan; reuses Ch6's digest machinery).

\section{Implementation}
\label{sec:spec-impl}
% The main/spec/ modules (parser, guard evaluator, executor fast path,
% state snapshot, commit log); executor as a mode branch of the agent
% loop, not a separate task; structured-output enforcement and
% malformed-script rejection (malformed rate reported as a result).
% Recovery from the commit log across resets.

\section{Evaluation}
\label{sec:spec-eval}
\subsection{Benefit at Matched Success}
% Calls/task, joules/task, the bimodal step-latency histogram, vs the
% four baselines (step ReAct; +prompt caching; ReWOO-style open loop;
% full-plan + naive equality guards).
\subsection{Depth versus Stochasticity}
% The core figure: acceptance length L and net benefit as environment
% noise/faults increase; the crossover where speculation stops paying.
\subsection{Guard Quality}
% The ROC-shaped trade-off: silent divergence vs false-reject across
% envelope tightness and guard density; the effect of forcing one
% observable post-condition per action.
\subsection{Disconnection Resilience}
% Task survival under 10/60/600 s outages, vs step ReAct and the
% cloud-orchestrated shim (type-iv baseline).

\section{Discussion}
\label{sec:spec-discussion}
% When to speculate and how deep (the practical decision rule); the
% relationship to LLM-generated offline rules (WireClaw) and to FSM
% compilation as the limit case (future work); threats to validity.
